\documentclass[10pt,a4paper]{article}
\usepackage[utf8]{inputenc}
\usepackage[T1]{fontenc}
\usepackage[brazilian]{babel}
\usepackage{lmodern}
\usepackage{amsmath,amssymb}
\usepackage[margin=1.7cm,top=1.7cm,bottom=1.7cm]{geometry}
\usepackage{booktabs,array}
\usepackage{xcolor}
\usepackage{fancyhdr}
\usepackage{listings}
\usepackage[hidelinks]{hyperref}

\definecolor{azul}{HTML}{183B56}
\definecolor{cinza}{RGB}{248,249,250}
\definecolor{cinzaBorda}{RGB}{210,215,220}
\definecolor{cinzaEscuro}{RGB}{80,80,80}

\pagestyle{fancy}
\fancyhf{}
\fancyhead[L]{\footnotesize\textbf{Projeto e Análise de Algoritmos (PAA)} -- Prof. Felipe A. Louza}
\fancyhead[R]{\footnotesize\textbf{Lista 7: O Problema da Seleção}}
\fancyfoot[L]{\footnotesize Arthur Pires Lopes (12221EEL001) -- UFU}
\fancyfoot[R]{\footnotesize Página \thepage\ de 4}
\setlength{\headheight}{12pt}
\setlength{\parindent}{0pt}
\setlength{\parskip}{2.5pt}

\newcommand{\res}[1]{\par\smallskip\centerline{\fbox{\bfseries\strut #1}}\par\smallskip}
\newcommand{\resmath}[1]{\[\boxed{#1}\]}
\newcommand{\topico}[1]{\smallskip\noindent\textbf{\color{azul}#1}\par}
\newcommand{\questao}[2]{\section*{\large\color{azul}Questão #1 -- #2}\vspace{-4pt}}

\lstset{
  basicstyle=\ttfamily\scriptsize,
  backgroundcolor=\color{cinza},
  frame=single,
  rulecolor=\color{cinzaBorda},
  columns=fullflexible,
  keepspaces=true,
  showstringspaces=false,
  numbers=left,
  numberstyle=\tiny\color{cinzaEscuro},
  numbersep=5pt,
  xleftmargin=12pt,
  framexleftmargin=10pt,
  literate={á}{{\'a}}1 {ã}{{\~a}}1 {ç}{{\c c}}1 {é}{{\'e}}1
           {í}{{\'i}}1 {ó}{{\'o}}1 {ú}{{\'u}}1 {Á}{{\'A}}1
           {É}{{\'E}}1 {Í}{{\'I}}1 {Ó}{{\'O}}1 {Ú}{{\'U}}1
           {à}{{\`a}}1 {õ}{{\~o}}1 {â}{{\^a}}1 {ê}{{\^e}}1
           {ô}{{\^o}}1 {<=}{{$\le$}}1 {>=}{{$\ge$}}1 {!=}{{$\neq$}}1
           {<-}{{$\leftarrow$}}1
}

\begin{document}

% ==============================================================================
% PÁGINA 1: CABEÇALHO + QUESTÃO 1
% ==============================================================================
\begin{center}
  \textbf{\large Universidade Federal de Uberlândia -- UFU \quad$|$\quad Faculdade de Engenharia Elétrica}\\[2pt]
  {\Large\bfseries\color{azul} Lista 7: O Problema da Seleção}\\[2pt]
  \textbf{Aluno:} Arthur Pires Lopes \quad$|$\quad \textbf{Matrícula:} 12221EEL001 \quad$|$\quad \textbf{Prof.:} Felipe A. Louza
\end{center}
\hrule
\vspace{4pt}

\questao{1}{Mínimo e Máximo Simultâneos com no Máximo $3\lfloor n/2 \rfloor$ Comparações}

\textbf{Enunciado:} \textit{Escreva o pseudocódigo de algoritmo que encontre o máximo e o mínimo em um vetor $A[1, n]$ realizando no máximo $3\lfloor n/2 \rfloor$ comparações.}

\topico{1. Estratégia Algorítmica (Processamento aos Pares)}
O método ingênuo compara cada elemento com o $\min$ e com o $\max$ corrente, consumindo $2n - 2$ comparações. 
Para reduzir o custo para no máximo $3\lfloor n/2 \rfloor$ comparações, processamos os elementos \textbf{aos pares}:
\begin{enumerate}
    \setlength{\itemsep}{0pt}
    \item Comparamos os dois elementos do par $(A[i], A[i+1])$ entre si ($1$ comparação) para definir o menor e o maior do par.
    \item Comparamos o menor do par com o $\min$ corrente ($1$ comparação).
    \item Comparamos o maior do par com o $\max$ corrente ($1$ comparação).
\end{enumerate}
Dessa forma, gastamos exatamente \textbf{3 comparações para cada 2 elementos} processados ($1{,}5$ comp./elem.).

\topico{2. Pseudocódigo}
\begin{lstlisting}
MIN-MAX(A, n)
 1  se n == 1 então retorne (A[1], A[1])
 2  se n é ímpar então
 3      min <- A[1];  max <- A[1];  inicio <- 2
 4  senão
 5      se A[1] <= A[2] então min <- A[1];  max <- A[2]
 6      senão min <- A[2];  max <- A[1]
 7      inicio <- 3
 8  para i <- inicio até n com passo 2 faça
 9      se A[i] <= A[i+1] então
10          se A[i] < min então min <- A[i]
11          se A[i+1] > max então max <- A[i+1]
12      senão
13          se A[i+1] < min então min <- A[i+1]
14          se A[i] > max então max <- A[i]
15  retorne (min, max)
\end{lstlisting}

\topico{3. Contagem Formal de Comparações de Chaves}
\begin{itemize}
    \setlength{\itemsep}{1pt}
    \item \textbf{Se $n$ é ímpar:} Inicialização direta com $A[1]$ custa $0$ comparações. Restam $n-1$ elementos, totalizando $\frac{n-1}{2} = \lfloor n/2 \rfloor$ pares. Cada par realiza 3 comparações no laço:
    \[
    C(n) = 0 + 3 \cdot \left(\frac{n-1}{2}\right) = 3\left\lfloor \frac{n}{2} \right\rfloor \text{ comparações.}
    \]
    \item \textbf{Se $n$ é par:} Inicialização compara $A[1]$ e $A[2]$ ($1$ comparação). Restam $n-2$ elementos, agrupados em $\frac{n-2}{2} = \lfloor n/2 \rfloor - 1$ pares. Cada par realiza 3 comparações:
    \[
    C(n) = 1 + 3 \cdot \left(\frac{n-2}{2}\right) = 1 + \frac{3n}{2} - 3 = \frac{3n}{2} - 2 = 3\left\lfloor \frac{n}{2} \right\rfloor - 2 \text{ comparações.}
    \]
\end{itemize}
Logo, em ambos os casos (par ou ímpar):
\resmath{C(n) \le 3\left\lfloor \frac{n}{2} \right\rfloor \quad \text{para todo } n \ge 1.}

\clearpage
% ==============================================================================
% PÁGINA 2: QUESTÃO 2
% ==============================================================================
\questao{2}{Execução Passo a Passo do Algoritmo \texttt{SELECTION}}

\textbf{Enunciado:} \textit{Mostre a execução passo-a-passo do algoritmo Selection para encontrar o 3º menor elemento para o seguinte vetor: $A = \langle 1, 8, 7, 6, 5, 3, 2, 10, 9 \rangle$. Considere que o pivot é escolhido sempre como o primeiro elemento do vetor.}

\topico{1. Mecânica do Algoritmo \texttt{SELECTION(A, p, r, i)}}
A função busca a $i$-ésima menor chave em $A[p \dots r]$. A partição fixa o pivô $x = A[p]$, rearranja o subvetor e posiciona o pivô no índice final $q$. O rank relativo do pivô no subvetor é $k = q - p + 1$:
\begin{itemize}
    \setlength{\itemsep}{0pt}
    \item Se $i = k$: o pivô $A[q]$ é o elemento procurado $\implies$ retorna $A[q]$;
    \item Se $i < k$: o elemento está na metade esquerda $\implies$ \texttt{SELECTION(A, p, q - 1, i)};
    \item Se $i > k$: o elemento está na metade direita $\implies$ \texttt{SELECTION(A, q + 1, r, i - k)}.
\end{itemize}

\topico{2. Rastreamento Passo a Passo ($n = 9$, busca do $3^\circ$ menor: $i = 3$)}

\begin{enumerate}
    \setlength{\itemsep}{2pt}
    \item \textbf{Chamada 1: \texttt{SELECTION(A, 1, 9, 3)}}
    \begin{itemize}
        \setlength{\itemsep}{0pt}
        \item Subvetor: $A[1 \dots 9] = \langle \mathbf{1}, 8, 7, 6, 5, 3, 2, 10, 9 \rangle$. Pivô: $x = A[1] = 1$.
        \item Varredura: nenhum elemento $\le 1$. Pivô permanece em $q = 1$. Vetor: $\langle \underline{1}, 8, 7, 6, 5, 3, 2, 10, 9 \rangle$.
        \item Rank: $k = 1 - 1 + 1 = 1$. Como $i = 3 > k = 1$: busca à direita com $i' = 3 - 1 = 2$.
    \end{itemize}

    \item \textbf{Chamada 2: \texttt{SELECTION(A, 2, 9, 2)}}
    \begin{itemize}
        \setlength{\itemsep}{0pt}
        \item Subvetor: $A[2 \dots 9] = \langle \mathbf{8}, 7, 6, 5, 3, 2, 10, 9 \rangle$. Pivô: $x = A[2] = 8$.
        \item Elementos $\le 8$: $\{7, 6, 5, 3, 2\}$ (5 chaves). O pivô $8$ troca com o último menor ($2$), parando em $q = 7$.
        \item Vetor: $\langle 1, \mathbf{2, 7, 6, 5, 3}, \underline{8}, \mathbf{10, 9} \rangle$. Rank: $k = 7 - 2 + 1 = 6$.
        \item Como $i = 2 < k = 6$: busca à esquerda com $i' = 2$.
    \end{itemize}

    \item \textbf{Chamada 3: \texttt{SELECTION(A, 2, 6, 2)}}
    \begin{itemize}
        \setlength{\itemsep}{0pt}
        \item Subvetor: $A[2 \dots 6] = \langle \mathbf{2}, 7, 6, 5, 3 \rangle$. Pivô: $x = A[2] = 2$.
        \item Todos os demais $\{7, 6, 5, 3\} > 2$. O pivô permanece em $q = 2$. Vetor: $\langle 1, \underline{2}, \mathbf{7, 6, 5, 3}, 8, 10, 9 \rangle$.
        \item Rank: $k = 2 - 2 + 1 = 1$. Como $i = 2 > k = 1$: busca à direita com $i' = 2 - 1 = 1$.
    \end{itemize}

    \item \textbf{Chamada 4: \texttt{SELECTION(A, 3, 6, 1)}}
    \begin{itemize}
        \setlength{\itemsep}{0pt}
        \item Subvetor: $A[3 \dots 6] = \langle \mathbf{7}, 6, 5, 3 \rangle$. Pivô: $x = A[3] = 7$.
        \item Menores ou iguais a $7$: $\{6, 5, 3\}$. Troca $A[3]$ com $A[6]$ ($7$ com $3$): pivô fixa-se em $q = 6$.
        \item Vetor: $\langle 1, 2, \mathbf{3, 6, 5}, \underline{7}, 8, 10, 9 \rangle$. Rank: $k = 6 - 3 + 1 = 4$.
        \item Como $i = 1 < k = 4$: busca à esquerda com $i' = 1$.
    \end{itemize}

    \item \textbf{Chamada 5: \texttt{SELECTION(A, 3, 5, 1)}}
    \begin{itemize}
        \setlength{\itemsep}{0pt}
        \item Subvetor: $A[3 \dots 5] = \langle \mathbf{3}, 6, 5 \rangle$. Pivô: $x = A[3] = 3$.
        \item Elementos $\{6, 5\} > 3$. O pivô permanece em $q = 3$. Vetor: $\langle 1, 2, \underline{3}, 6, 5, 7, 8, 10, 9 \rangle$.
        \item Rank: $k = 3 - 3 + 1 = 1$. Como $i = 1 == k = 1$: \textbf{condição de parada!}
        \item O algoritmo retorna o elemento no pivô: $A[3] = 3$.
    \end{itemize}
\end{enumerate}

\topico{3. Tabela Resumo das Chamadas}
\begin{center}
\footnotesize
\setlength{\tabcolsep}{5pt}
\begin{tabular}{ccccccc}
\toprule
\textbf{Chamada} & \textbf{Subvetor} & \textbf{Pivô ($x$)} & \textbf{Posição ($q$)} & \textbf{Rank ($k$)} & \textbf{Ordem ($i$)} & \textbf{Ação Realizada} \\ \midrule
1 & $A[1 \dots 9]$ & $A[1]=1$ & 1 & 1 & 3 & $i > k \implies$ Recurso à direita ($i \leftarrow 2$) \\
2 & $A[2 \dots 9]$ & $A[2]=8$ & 7 & 6 & 2 & $i < k \implies$ Recurso à esquerda ($i \leftarrow 2$) \\
3 & $A[2 \dots 6]$ & $A[2]=2$ & 2 & 1 & 2 & $i > k \implies$ Recurso à direita ($i \leftarrow 1$) \\
4 & $A[3 \dots 6]$ & $A[3]=7$ & 6 & 4 & 1 & $i < k \implies$ Recurso à esquerda ($i \leftarrow 1$) \\
5 & $A[3 \dots 5]$ & $A[3]=3$ & 3 & 1 & 1 & \textbf{$i = k \implies$ Retorna $A[3] = 3$} \\
\bottomrule
\end{tabular}
\end{center}

\res{O $3^\circ$ menor elemento retornado é: \boxed{3}}

\clearpage
% ==============================================================================
% PÁGINA 3: QUESTÃO 3
% ==============================================================================
\questao{3}{Comparação de Métodos para Listar os $i$ Maiores Elementos em Ordem Crescente}

\textbf{Enunciado:} \textit{Dado um vetor $A[1, n]$, queremos imprimir em ordem crescente os $i$ maiores elementos de $A$. Compare a complexidade dos seguintes métodos em função de $n$ e $i$:}
\textit{(a) Ordene o vetor e liste os $i$ maiores elementos.}
\textit{(b) Construa uma fila de prioridade (max-heap) e chame a rotina para extrair o máximo $i$ vezes.}
\textit{(c) Use um algoritmo de seleção para encontrar o $i$-ésimo maior elemento, particione o vetor em torno dele e ordene os $i$ maiores elementos.}

\topico{1. Método (a): Ordenação Total de $A[1 \dots n]$}
\begin{itemize}
    \setlength{\itemsep}{1pt}
    \item \textbf{Ordenação de todo o vetor:} Algoritmo ótimo baseado em comparações (MergeSort/HeapSort) custa $\Theta(n \lg n)$.
    \item \textbf{Listagem dos $i$ maiores:} Em vetor ordenado crescente, os $i$ maiores estão nas posições $A[n - i + 1 \dots n]$. Percorrer e imprimir essas posições consome tempo $\Theta(i)$.
    \item \textbf{Complexidade Total:} $T_a(n, i) = \Theta(n \lg n) + \Theta(i) = \Theta(n \lg n)$ (visto que $1 \le i \le n$).
\end{itemize}

\topico{2. Método (b): Max-Heap com $i$ Extrações de Máximo}
\begin{itemize}
    \setlength{\itemsep}{1pt}
    \item \textbf{Construção do Max-Heap (\texttt{BUILD-MAX-HEAP}):} Constrói o heap linearmente em tempo $\Theta(n)$.
    \item \textbf{Extrações de máximo ($i$ vezes):} Cada \texttt{EXTRACT-MAX} remove a raiz e repara a estrutura via \texttt{MAX-HEAPIFY} em tempo $O(\lg n)$. O custo acumulado de $i$ extrações é:
    \[
    \sum_{j=0}^{i-1} O(\lg(n - j)) = \Theta(i \lg n).
    \]
    \item \textbf{Ajuste para ordem crescente:} As extrações retornam elementos decrescentes; ordená-los crescentemente (armazenando de trás para frente em vetor auxiliar de tamanho $i$) leva $\Theta(i)$.
    \item \textbf{Complexidade Total:} $T_b(n, i) = \Theta(n) + \Theta(i \lg n) + \Theta(i) = \Theta(n + i \lg n)$.
\end{itemize}

\topico{3. Método (c): Seleção Linear, Partição e Ordenação Parcial de Tamanho $i$}
\begin{itemize}
    \setlength{\itemsep}{1pt}
    \item \textbf{Seleção da estatística de ordem:} O $i$-ésimo maior elemento é o $(n - i + 1)$-ésimo menor de $A$. Utilizando algoritmo determinístico linear (\texttt{SELECT} mediana das medianas, Cormen Cap. 9), o custo é $\Theta(n)$.
    \item \textbf{Particionamento:} Particionar $A$ em torno dele posiciona os $i$ maiores no bloco final $A[n - i + 1 \dots n]$ em $\Theta(n)$.
    \item \textbf{Ordenação parcial:} Ordenamos apenas o subvetor de tamanho $i$ em ordem crescente (MergeSort), custando $\Theta(i \lg i)$.
    \item \textbf{Listagem:} Imprimir os $i$ elementos já ordenados custa $\Theta(i)$.
    \item \textbf{Complexidade Total:} $T_c(n, i) = \Theta(n) + \Theta(n) + \Theta(i \lg i) + \Theta(i) = \Theta(n + i \lg i)$.
\end{itemize}

\topico{4. Quadro Comparativo e Análise Assintótica}
\begin{center}
\small
\setlength{\tabcolsep}{5pt}
\begin{tabular}{p{1.5cm}p{8.5cm}p{3.8cm}}
\toprule
\textbf{Método} & \textbf{Etapas Fundamentais} & \textbf{Complexidade Total} \\ \midrule
\textbf{(a)} & Ordenação total de $A$ + listagem dos $i$ maiores & $\Theta(n \lg n)$ \\
\textbf{(b)} & \texttt{BUILD-MAX-HEAP} + $i \times$ \texttt{EXTRACT-MAX} + inversão & $\Theta(n + i \lg n)$ \\
\textbf{(c)} & \texttt{SELECT} linear + partição + ordenação de $i$ elementos & $\Theta(n + i \lg i)$ \\
\bottomrule
\end{tabular}
\end{center}

Como $1 \le i \le n$, vale sempre que $\lg i \le \lg n$. Desse modo:
\[
\Theta(n + i \lg i) \le \Theta(n + i \lg n) \le \Theta(n \lg n).
\]
\begin{itemize}
    \setlength{\itemsep}{1pt}
    \item Se $i = O(1)$ ou $i = O(n/\lg n)$: Métodos (b) e (c) operam em tempo linear $\Theta(n)$, muito superiores a (a).
    \item Se $i = \Theta(n)$: Todos têm ordem $\Theta(n \lg n)$, mas (c) possui menor constante oculta ($\alpha \lg(\alpha n) < \lg n$).
\end{itemize}
\resmath{T_c(n, i) = \Theta(n + i \lg i) \quad \text{é assintoticamente o mais eficiente para qualquer } i.}

\clearpage
% ==============================================================================
% PÁGINA 4: QUESTÕES 4 E 5
% ==============================================================================
\questao{4}{QuickSort Modificado com Tempo $\Theta(n \lg n)$ no Pior Caso}

\textbf{Enunciado:} \textit{Mostre como modificar QuickSort de modo que tenha complexidade de tempo $\Theta(n \lg n)$ no pior caso.}

\topico{1. Diagnóstico e Modificação}
O \texttt{QuickSort} tradicional degenera para $\Theta(n^2)$ quando ocorrem partições assimétricas ($0$ e $n-1$).
Para garantir tempo ótimo $\Theta(n \lg n)$ no pior caso, selecionamos como pivô a \textbf{mediana exata} do subvetor usando a sub-rotina determinística \texttt{SELECT} (mediana das medianas / BFPRT), que roda em tempo linear $\Theta(m)$ no pior caso.

\topico{2. Pseudocódigo e Análise da Recorrência}
\begin{minipage}[t]{0.48\textwidth}
\begin{lstlisting}
QUICKSORT-OTIMO(A, p, r)
 1  se p < r então
 2      m <- r - p + 1
 3      k <- piso((m + 1) / 2)
 4      med <- SELECT(A, p, r, k) // O(m)
 5      idx <- indice de 'med' em [p, r]
 6      troca A[idx] com A[r]
 7      q <- PARTITION(A, p, r)  // Theta(m)
 8      QUICKSORT-OTIMO(A, p, q - 1)
 9      QUICKSORT-OTIMO(A, q + 1, r)
\end{lstlisting}
\end{minipage}
\hfill
\begin{minipage}[t]{0.48\textwidth}
\textbf{Demonstração via Teorema Mestre:}
Com a mediana como pivô, cada subproblema tem no máximo $\lfloor m/2 \rfloor \le n/2$ elementos.
Trabalho não recursivo: \texttt{SELECT} $O(n)$ $+$ troca $O(n)$ $+$ \texttt{PARTITION} $\Theta(n) \implies cn$.

Recorrência de pior caso:
\[
T(n) \le 2\,T(n/2) + cn.
\]
Com $a = 2, b = 2 \implies n^{\log_2 2} = n^1 = n$. Como $f(n) = cn \in \Theta(n)$, estamos no \textbf{Caso 2 do Teorema Mestre}:
\[
T(n) = \Theta(n \lg n).
\]
\end{minipage}

Como qualquer ordenação por comparações exige $\Omega(n \lg n)$ comparações, a cota é estrita:
\resmath{T(n) \in \Theta(n \lg n) \quad \text{no pior caso.}}

\vspace{4pt}
\hrule
\vspace{4pt}

\questao{5}{Algoritmo Linear Simples de Seleção Usando Caixa-Preta de Mediana}

\textbf{Enunciado:} \textit{Suponha que você tenha uma subrotina do tipo ``caixa-preta'' que determina a mediana em tempo linear (no pior caso). Descreva um algoritmo linear simples que resolve o problema da seleção para todo $i$.}

\topico{1. Descrição e Pseudocódigo}
Seja \texttt{MEDIANA(A, p, r)} a caixa-preta que determina a mediana em tempo $c_1 m$ ($m = r - p + 1$). Usamos essa mediana como pivô perfeito para particionar o vetor. Diferente do QuickSort, a seleção faz \textbf{apenas uma chamada recursiva} no subvetor que contém a estatística procurada:

\begin{lstlisting}
SELECAO-LINEAR(A, p, r, i)
 1  se p == r então retorne A[p]
 2  med <- MEDIANA(A, p, r)                 // Caixa-preta: custo O(m) no pior caso
 3  idx <- indice de 'med' em [p, r];  troca A[idx] com A[r]
 4  q <- PARTITION(A, p, r)                 // Particionamento linear Theta(m)
 5  k <- q - p + 1                          // Posto relativo do pivo
 6  se i == k então retorne A[q]
 7  senão se i < k então retorne SELECAO-LINEAR(A, p, q - 1, i)
 8  senão retorne SELECAO-LINEAR(A, q + 1, r, i - k)
\end{lstlisting}

\topico{2. Demonstração da Complexidade Linear $\Theta(n)$}
A mediana divide o vetor em dois lados com no máximo $\lfloor n/2 \rfloor \le n/2$ chaves. Como apenas uma metade é explorada e o custo não recursivo por etapa é $d \cdot n$, a relação de recorrência no pior caso é:
\[
T(n) \le T(n/2) + dn.
\]
Expandindo pela soma da série geométrica de razão $1/2$:
\[
T(n) \le dn \sum_{j=0}^{\infty} \left(\frac{1}{2}\right)^j = dn \cdot \left(\frac{1}{1 - 1/2}\right) = 2dn = O(n).
\]
Como qualquer algoritmo de seleção precisa examinar os $n$ elementos no pior caso ($T(n) \in \Omega(n)$):
\resmath{T(n) \in \Theta(n) \quad \text{no pior caso, para qualquer } i.}

\end{document}
